% !TeX program = xelatex

\documentclass[aspectratio=169]{beamer}

% ============================================================
% 1. Fonts
% ============================================================
\usepackage{fontspec}

\setmainfont{Microsoft JhengHei}
\setsansfont{Microsoft JhengHei}
\setmonofont{JetBrains Mono}

\XeTeXlinebreaklocale "zh"


% ============================================================
% 2. Mathematics, enumerate
% ============================================================
\usepackage{amsmath}
\usepackage{amssymb}
\usepackage{enumitem}

\setlistdepth{6}
\renewlist{itemize}{itemize}{9}
\renewlist{enumerate}{enumerate}{9}
\setlist[itemize,1]{label=\textbullet}
\setlist[itemize,2]{label=$\circ$}
\setlist[itemize,3]{label=$\ast$}
\setlist[itemize,4]{label=$\cdot$}
\setlist[itemize,5]{label=--}
\setlist[itemize,6]{label=--}
\setlist[enumerate,1]{label=\arabic*.}
\setlist[enumerate,2]{label=\arabic*.}
\setlist[enumerate,3]{label=\arabic*.}
\setlist[enumerate,4]{label=\arabic*.}


% ============================================================
% 3. Images / Tables / Boxes
% ============================================================
\usepackage{graphicx}
\usepackage{booktabs}
\usepackage[most]{tcolorbox}


% ============================================================
% 4. Colors
% ============================================================
\definecolor{myRed}{RGB}{145,90,90}
\definecolor{myGreen}{RGB}{85,115,100}
\definecolor{myBlue}{RGB}{70,100,130}
\definecolor{myOrange}{RGB}{155,115,80}
\definecolor{myPurple}{RGB}{115,100,130}

\definecolor{myGray}{RGB}{242,242,238}
\definecolor{myDarkgray}{RGB}{90,90,88}
\definecolor{myBlack}{RGB}{45,45,43}


% ============================================================
% 5. Beamer Colors
% ============================================================
\setbeamercolor{structure}   {fg=myBlue}
\setbeamercolor{title}       {fg=white, bg=myBlue}
\setbeamercolor{frametitle}  {fg=white, bg=myBlue}
\setbeamercolor{block title} {fg=white, bg=myBlue}
\setbeamercolor{block body}  {fg=myBlack, bg=myGray}
\setbeamercolor{alerted text}{fg=myRed}
\setbeamercolor{example text}{fg=myGreen}


% ============================================================
% 6. Theme
% ============================================================
\usetheme{Madrid}
\usecolortheme{default}

\setbeamertemplate{navigation symbols}{}


% ============================================================
% 7. Footer
% ============================================================
\setbeamertemplate{footline}{
    \leavevmode
    \hbox{
        % Left: Author
        \begin{beamercolorbox}[
            wd=.33\paperwidth,
            ht=2.5ex,
            dp=1ex,
            center
        ]{author in head/foot}
            \insertdate
        \end{beamercolorbox}

        % Center: Title
        \begin{beamercolorbox}[
            wd=.33\paperwidth,
            ht=2.5ex,
            dp=1ex,
            center
        ]{title in head/foot}
            \insertshorttitle
        \end{beamercolorbox}

        % Right: Page Number
        \begin{beamercolorbox}[
            wd=.33\paperwidth,
            ht=2.5ex,
            dp=1ex,
            center
        ]{date in head/foot}
            \insertframenumber{} / \inserttotalframenumber
        \end{beamercolorbox}
    }
}


% ============================================================
% 8. Presentation Information
% ============================================================
\title[第三次雙周會]{畢業專題 文獻探討}
\subtitle{Agentic AI Security}
\author{Syu}
\date{October 19, 2026}



% ============================================================
% Helpers (table spacing)
% ============================================================
\newcommand{\tabsetup}{\renewcommand{\arraystretch}{1.18}}


% ========================= DocStart =========================
\begin{document}

\newcommand{\paperref}[1]{\vfill {\tiny\color{myDarkgray} arXiv: #1}}

\begin{frame}
    \titlepage
\end{frame}

\begin{frame}{Index}
    \tableofcontents
\end{frame}

\section{Review}
\begin{frame}{Review: Agentic AI Security}
    \begin{enumerate}
        \item \textbf{Agentic AI Security (2510.23883)}
            \begin{itemize}
                \item Threats: prompt injection, tool abuse, MCP, and multi-agent attacks.
                \item Defenses: sandboxing, access control, and runtime monitoring.
            \end{itemize}
        \item \textbf{Agentic AI and Cybersecurity (2601.05293)}
            \begin{itemize}
                \item Dual-use agents: autonomous defense and automated exploitation.
            \end{itemize}
        \item \textbf{Agent-SafetyBench (2412.14470)}
            \begin{itemize}
                \item Interactive safety evaluation of agent actions and attack scenarios.
            \end{itemize}
    \end{enumerate}
    \begin{block}{Key Takeaway}
        Agent security requires enforceable controls beyond model-level alignment.
    \end{block}
\end{frame}

\begin{frame}{Review: Streamo}
    \begin{block}{Streaming Video Instruction Tuning}
        \begin{itemize}
            \item \textbf{Problem:} Offline video models and multi-pass pipelines increase response latency.
            \item \textbf{Method:} One-pass inference with three state tokens:
                \begin{itemize}
                    \item \texttt{<Silence>}: Continue processing frames.
                    \item \texttt{<Standby>}: Wait for an ongoing event.
                    \item \texttt{<Response>}: Generate an answer.
                \end{itemize}
            \item \textbf{Training:} Streamo-Instruct-465K and focal-weighted LM loss.
        \end{itemize}
    \end{block}
    \begin{block}{Key Takeaway}
        Streamo prioritizes real-time interaction; Video Streaming Thinking emphasizes longer-term reasoning.
    \end{block}
\end{frame}

\begin{frame}{Review: StreamEQA}
    \begin{block}{Streaming Embodied Question Answering}
        \begin{itemize}
            \item \textbf{Goal:} Evaluate online, first-person video understanding without access to future frames.
            \item \textbf{Embodied dimensions:}
                \begin{itemize}
                    \item Perception, Interaction, Planning.
                \end{itemize}
            \item \textbf{Streaming dimensions:}
                \begin{itemize}
                    \item Backward Reasoning, Realtime Understanding, Forward Prediction.
                \end{itemize}
            \item \textbf{Method:} Extract video metadata, construct QA pairs, and benchmark multimodal models.
        \end{itemize}
    \end{block}
    \begin{block}{Connection}
        Real-time agent decisions also need safety checks before external actions execute.
    \end{block}
\end{frame}

\section{Pre-Action Authorization}
\begin{frame}{Pre-Action Authorization (arXiv 2603.20953)}
    \footnotesize
    \textbf{Before the Tool Call: Deterministic Pre-Action Authorization for Autonomous AI Agents}

    {\small Uchi Uchibeke\\APort Technologies Inc., Toronto, Canada}
    \begin{block}{Abstract}
        Autonomous agents can invoke powerful tools without a standardized, deterministic authorization boundary.
    \end{block}
    \begin{itemize}
        \item \textbf{OAP:} Intercepts each tool call before execution.
        \item Evaluates declarative policies against agent identity, capabilities, and call arguments.
        \item Produces cryptographically signed audit records.
        \item \textbf{Result:} 53 ms median authorization latency; 0 successful attacks in 879 restrictive-policy attempts.
    \end{itemize}
    \paperref{2603.20953v1}
\end{frame}

\begin{frame}{Authorization Model (arXiv 2603.20953)}
    \begin{block}{Deterministic Authorization}
        \[\operatorname{authorize}(T,P,\Pi)\longrightarrow(D,L)\]
        \begin{itemize}
            \item $T$: Tool call and arguments; $P$: Agent passport and capabilities.
            \item $\Pi$: Declarative policies; $L$: Signed audit log.
            \item $D\in\{\text{ALLOW},\text{DENY},\text{ESCALATE}\}$.
        \end{itemize}
    \end{block}
    \begin{enumerate}
        \item Validate identity, capabilities, and policy constraints.
        \item Reject unauthorized arguments or insufficient assurance.
        \item Issue a decision before execution and record the result.
    \end{enumerate}
    \begin{alertblock}{Boundary}
        Deterministic authorization complements, but does not replace, sandbox isolation.
    \end{alertblock}
    \paperref{2603.20953v1}
\end{frame}

\begin{frame}{Adversarial Evaluation (arXiv 2603.20953)}
    \begin{block}{Results}
        \begin{itemize}
            \item \textbf{53 ms} median decision latency ($N=1{,}000$).
            \item \textbf{74.6\%} model-level social-engineering success under permissive policies.
            \item \textbf{0/879} attack successes under restrictive OAP policies.
        \end{itemize}
    \end{block}
    \begin{alertblock}{Limitations}
        \begin{itemize}
            \item Static policy design; cross-agent delegation is not fully formalized.
            \item Individually authorized actions may still compose into unsafe workflows.
            \item Human approval (ESCALATE) is not implemented in the reference system.
        \end{itemize}
    \end{alertblock}
    \paperref{2603.20953v1}
\end{frame}

\section{MCP Security Framework}
\begin{frame}{MCP Threat Modeling (arXiv 2604.05969)}
    \footnotesize
    \textbf{A Formal Security Framework for MCP-Based AI Agents: Threat Taxonomy, Verification Models, and Defense Mechanisms}

    {\small Nirajan Acharya (University of the Cumberlands)\\Gaurav Kumar Gupta (Youngstown State University)}
    \begin{block}{Abstract}
        MCP tool ecosystems introduce new trust boundaries and fragmented security controls.
    \end{block}
    \begin{itemize}
        \item \textbf{Taxonomy:} 7 threat categories, 23 attack vectors, 4 attack surfaces.
        \item \textbf{Method:} Labeled transition systems and formal security properties.
        \item \textbf{Analysis:} Comparison of 12 existing defense mechanisms.
        \item \textbf{Proposal:} Four-layer defense-in-depth architecture with 91\% theoretical coverage.
    \end{itemize}
    \paperref{2604.05969v1}
\end{frame}

\begin{frame}{Formal Verification Model (arXiv 2604.05969)}
    \begin{block}{MCP Model and Labeled Transition System}
        \[M=(A,S,T,R,C,\ell),\quad LTS=(\Sigma,\Sigma_0,\Lambda,\rightarrow)\]
        \small
        $A$: Agents; $S$: Servers; $T$: Tools; $R$: Resources;\\
        $C$: Connections; $\ell$: Security labels; $\Sigma$: System states.
    \end{block}
    \begin{block}{Verified Properties}
        \begin{itemize}
            \item \textbf{Tool integrity:} Prevent post-approval tool modifications.
            \item \textbf{Data confinement:} Restrict cross-boundary information flows.
            \item \textbf{Privilege boundedness:} Enforce capability limits.
            \item \textbf{Context isolation:} Block unauthorized context sharing.
        \end{itemize}
    \end{block}
    \paperref{2604.05969v1}
\end{frame}

\begin{frame}{Defense-in-Depth (arXiv 2604.05969)}
    \small
    \begin{block}{Proposed Layers}
        \begin{enumerate}
            \item Capability-based access control.
            \item Cryptographic tool attestation.
            \item Information-flow tracking.
            \item Runtime policy enforcement.
        \end{enumerate}
    \end{block}
    \begin{block}{Coverage Analysis}
        \begin{itemize}
            \item No individual defense covers more than \textbf{34\%} of the identified threats.
            \item Integrated architecture: \textbf{91\% theoretical threat coverage}.
        \end{itemize}
    \end{block}
    \small\textbf{Note:} Coverage is theoretical, not a measured block rate.
    \paperref{2604.05969v1}
\end{frame}

\section{Learned Capability Governance}
\begin{frame}{Adaptive Least Privilege (arXiv 2604.11839)}
    \footnotesize
    \textbf{Beyond Static Sandboxing: Learned Capability Governance for Autonomous AI Agents}

    {\small Bronislav Sidik, Lior Rokach\\Ben-Gurion University of the Negev, Institute for Applied AI Research}
    \begin{block}{Abstract}
        Agent runtimes expose excessive tools even for simple tasks, expanding the attack surface.
    \end{block}
    \begin{itemize}
        \item \textbf{AgentWarden:} Task-aware capability scoping and pre-execution safety routing.
        \item \textbf{Learning:} PPO optimizes tool availability using session audit data.
        \item \textbf{Metric:}
            \[\mathrm{SER}=\frac{|\text{Tools invoked}|}{|\text{Tools exposed}|}\]
        \item \textbf{Result:} 73\% fewer exposed tools for summarization tasks.
    \end{itemize}
    \paperref{2604.11839v2}
\end{frame}

\begin{frame}{MDP and PPO Policy (arXiv 2604.11839)}
    \begin{block}{Markov Decision Process}
        \begin{itemize}
            \item \textbf{State:} Task category and normalized trust level (7 dimensions).
            \item \textbf{Action:} Binary tool-access mask, $a\in\{0,1\}^{17}$.
            \item \textbf{Policy:} PPO with a two-layer MLP (64 units per layer).
        \end{itemize}
        \[R(s,a)=\alpha U_{\rm accuracy}-\beta C_{\rm economy}-\gamma P_{\rm safety}\]
        \small $(\alpha,\beta,\gamma)=(1,0.3,5)$.
    \end{block}
    \begin{block}{Runtime Enforcement}
        Capability Governor limits visible tools; Safety Router checks each proposed call.
    \end{block}
    \paperref{2604.11839v2}
\end{frame}

\begin{frame}{Evaluation and Failure Modes (arXiv 2604.11839)}
    \begin{block}{OpenClaw Experiments}
        \begin{itemize}
            \item \textbf{500 tasks:} 400 benign and 100 adversarial.
            \item \textbf{73\%} reduction in exposed tools for summarization.
            \item \textbf{26.2\%} of 447 intercepted tool calls blocked.
            \item \textbf{92\%} of adversarial tasks neutralized by the combined stack.
        \end{itemize}
    \end{block}
    \begin{alertblock}{Key Limitations}
        \begin{itemize}
            \item PPO sometimes collapsed to a single-tool policy during real evaluation.
            \item The system then fell back to YAML-based policies.
            \item Safety Router training data were small and imbalanced.
        \end{itemize}
    \end{alertblock}
    \paperref{2604.11839v2}
\end{frame}

\section{System-Wide Agent Governance}
\begin{frame}{Cross-Agent Action Governance (arXiv 2609.00015)}
    \footnotesize
    \textbf{OpenAgentFlow: Enabling System-Wide Safety Boundaries for Heterogeneous AI Agent Fleets}

    {\scriptsize Dongsheng Chen (Southern University of Science and Technology)\\Xiangyu Zhao (City University of Hong Kong)\\Xin Yao (Lingnan University); Xuetao Wei (Southern University of Science and Technology)}
    \begin{block}{Abstract}
        Individually valid actions can form unsafe workflows across multiple agents and applications.
    \end{block}
    \begin{itemize}
        \item Normalize GUI, API, and tool actions into shared \textbf{AgentEvents}.
        \item Enforce policy at a shared \textbf{action-commit boundary}.
        \item Maintain provenance, session state, and audit records outside individual agents.
        \item \textbf{Result:} 97.62\% accuracy on 1,220 AgentDojo-Traj cases.
    \end{itemize}
    \paperref{2609.00015v2}
\end{frame}

\begin{frame}{Shared Policy Enforcement (arXiv 2609.00015)}
    \begin{block}{Control Plane / Action Plane}
        \begin{itemize}
            \item \textbf{Control plane:} Policies, provenance, session state, audit records.
            \item \textbf{Action plane:} Shared Policy Enforcement Point (PEP).
            \item \textbf{Action-commit boundary:} Inspect an action before side effects occur.
        \end{itemize}
    \end{block}
    \begin{block}{Policy Decision Stages}
        \begin{enumerate}
            \item T1: Structured policy rules.
            \item T2: Payload and provenance checks.
            \item T3: Semantic assessment when rules are insufficient.
            \item T4: Final adjudication and escalation.
        \end{enumerate}
    \end{block}
    \paperref{2609.00015v2}
\end{frame}

\begin{frame}{System Evaluation (arXiv 2609.00015)}
    \begin{block}{Results}
        \begin{itemize}
            \item \textbf{300-case controlled suite:} 94.00\% accuracy; 95.35\% attack-block rate.
            \item \textbf{1,220-case AgentDojo-Traj:} 97.62\% accuracy.
                \begin{itemize}
                    \item 96.59\% unsafe-action recall.
                    \item 1.96\% false intervention on safe actions.
                \end{itemize}
            \item Demonstrated policy updates and enforcement on Android action paths.
        \end{itemize}
    \end{block}
    \begin{alertblock}{Key Challenge}
        Effective protection depends on complete event capture and accurate cross-step provenance.
    \end{alertblock}
    \paperref{2609.00015v2}
\end{frame}

\end{document}
